% =========================================================
% Ordered Sets
% =========================================================

\section{Ordered Sets}\label{sec:ordered-sets}

% ---------------------------------------------------------
% TOOLKIT
% ---------------------------------------------------------

\begin{definitionbox}{Definition (Ordered Set)}
\begin{definition}[Ordered Set]\label{def:ordered-set}
An \emph{ordered set} (or \emph{partially ordered set}, \emph{poset}) is a
pair $(A,\leq)$ where $\leq$ is a partial order on $A$: a relation that is
reflexive, antisymmetric, and transitive.
\end{definition}
\end{definitionbox}

\begin{dependencies}
\begin{itemize}
  \item \hyperref[def:partial-order]{Partial Order}
\end{itemize}
\end{dependencies}

\begin{remark*}[English reading]
An ordered set is a set equipped with a specified way to compare elements,
but without requiring all pairs to be comparable.
\end{remark*}

\begin{definitionbox}{Definition (Strict Order)}
\begin{definition}[Strict Order]\label{def:strict-order}
Let $(A,\leq)$ be an ordered set. The \emph{strict order} $<$ is defined by:
\[
a < b \;\;\Longleftrightarrow\;\; (a \leq b \land a \neq b).
\]
\end{definition}
\end{definitionbox}

\begin{remark*}[Strict and non-strict are equivalent]
The relation $<$ is irreflexive and transitive. Conversely, given any strict
partial order $<$, one recovers a non-strict order by setting
$a \leq b \iff (a < b \lor a=b)$. The two formulations carry exactly the
same information.
\end{remark*}

\begin{definitionbox}{Definition (Comparable and Incomparable Elements)}
\begin{definition}[Comparable and Incomparable Elements]\label{def:comparable}
In an ordered set $(A,\leq)$, elements $a,b \in A$ are \emph{comparable} if
$a \leq b$ or $b \leq a$; they are \emph{incomparable} if neither holds.
\end{definition}
\end{definitionbox}

\NoLocalDependencies

\begin{remark*}[Standard quantified statement]
$a,b$ comparable $\;\leftrightarrow\; a \leq b \vee b \leq a$. \quad
$a,b$ incomparable $\;\leftrightarrow\; \neg(a \leq b) \wedge \neg(b \leq a)$.
\end{remark*}

\begin{remark*}[Incomparable elements]
Incomparable elements can only occur in partial orders. In a total order,
every pair is comparable by definition.
\end{remark*}

\begin{definitionbox}{Definition (Total / Linear Order)}
\begin{definition}[Total / Linear Order]\label{def:linear-order}
An ordered set $(A,\leq)$ is a \emph{total order} (or \emph{linear order})
if every pair of elements is comparable:
\[
\forall a,b \in A,\; a \leq b \lor b \leq a.
\]
\end{definition}
\end{definitionbox}

\begin{remark*}[Examples]
$(\mathbb{R},\leq)$ is a total order. $(\mathcal{P}(A), \subseteq)$ is a
partial order that is generally not total.
\end{remark*}

% ---------------------------------------------------------
% Bounds and extremal elements
% ---------------------------------------------------------
\begin{definitionbox}{Definition (Upper and Lower Bounds)}
\begin{definition}[Upper and Lower Bounds]\label{def:bounds}
Let $(A,\leq)$ be an ordered set and $S \subseteq A$.

An element $u \in A$ is an \emph{upper bound} of $S$ if $\forall s \in S,\; s \leq u$.

An element $\ell \in A$ is a \emph{lower bound} of $S$ if $\forall s \in S,\; \ell \leq s$.
\end{definition}
\end{definitionbox}

\begin{remark*}[Consequence]
Bounds need not exist and need not be unique. They need not lie in $S$ itself.
Bounds are central to the definition of completeness for ordered fields.
\end{remark*}

\begin{propositionbox}{Definition (Minimal, Maximal, Least, Greatest Elements)}
\label{def:min-max}
Let $S \subseteq A$ in an ordered set $(A,\leq)$.

$m \in S$ is a \emph{minimal element} of $S$ if $\nexists s \in S$ with $s < m$.

$M \in S$ is a \emph{maximal element} of $S$ if $\nexists s \in S$ with $M < s$.

\smallskip
\label{def:least-greatest}
$\ell \in S$ is the \emph{least element} if $\forall s \in S,\; \ell \leq s$.

$g \in S$ is the \emph{greatest element} if $\forall s \in S,\; s \leq g$.
\end{propositionbox}

\begin{remark*}[Minimal vs.\ least]
A \emph{least} element is below every element of $S$; it must be unique if it
exists. A \emph{minimal} element merely has nothing below it within $S$; there
may be many. Every least element is minimal, but not conversely.
\end{remark*}

% ---------------------------------------------------------
% Order maps and isomorphisms
% ---------------------------------------------------------
\begin{propositionbox}{Definition (Order-Preserving Map and Isomorphism)}
\label{def:order-pres}
Let $(M,\leq)$ and $(M',\leq')$ be partially ordered sets.

A function $f : M \to M'$ is \emph{order-preserving} (or \emph{monotone}) if
\[
\forall a,b \in M,\quad a \leq b \;\Longrightarrow\; f(a) \leq' f(b).
\]

$f$ is an \emph{order isomorphism}\label{def:order-iso} if it is bijective and
\[
\forall a,b \in M,\quad a \leq b \;\Longleftrightarrow\; f(a) \leq' f(b).
\]
\end{propositionbox}

\begin{remark*}[Order-isomorphic posets]
An order isomorphism preserves and reflects the order structure exactly,
including comparability, minimal and maximal elements, and bounds. Two posets
are \emph{order-isomorphic} if such a map exists; they are structurally
identical from the order-theoretic viewpoint.
\end{remark*}

% ---------------------------------------------------------
% Well-order
% ---------------------------------------------------------
\begin{definitionbox}{Definition (Well-Ordered Set)}
\begin{definition}[Well-Ordered Set]\label{def:well-order}
An ordered set $(A,<)$ is \emph{well-ordered} if every nonempty subset
$S \subseteq A$ has a least element:
\[
\forall S \subseteq A,\;
\bigl(S \neq \varnothing \;\Rightarrow\; \exists m \in S \text{ s.t.\ }
m \leq s \;\forall s \in S\bigr).
\]
\end{definition}
\end{definitionbox}

\begin{remark*}[Well-order implies total order]
Every well-ordered set is totally ordered. The well-ordering condition is
strictly stronger: it requires a least element in every nonempty subset,
not just in $A$ as a whole.
\end{remark*}

\begin{example*}[Well-order examples]
$(\mathbb{N},\leq)$ is well-ordered. $(\mathbb{Z},\leq)$ is not: the set
$\{\dots,-3,-2,-1\}$ has no least element. $(\mathbb{R},\leq)$ is not:
$(0,1)$ has no least element.
\end{example*}

\begin{remark*}[Connection to ordinal numbers]
An \emph{ordinal number} is defined as an equivalence class of well-ordered
sets under order isomorphism: it measures the \emph{order type} of a
well-ordered set, not merely its cardinality.
\end{remark*}

% ---------------------------------------------------------
% Chains, antichains, initial segments
% ---------------------------------------------------------
\begin{definitionbox}{Definition (Chain and Antichain)}
\begin{definition}[Chain and Antichain]\label{def:chain}
A subset $C \subseteq A$ of a poset $(A,\leq)$ is a \emph{chain} if every
pair of elements in $C$ is comparable. It is an \emph{antichain} if no two
distinct elements of $C$ are comparable.
\end{definition}
\end{definitionbox}

\NoLocalDependencies

\begin{remark*}[Standard quantified statement]
$C$ is a chain $\;\leftrightarrow\; \forall a,b \in C,\; a \leq b \vee b \leq a$.\\
$C$ is an antichain $\;\leftrightarrow\; \forall a,b \in C,\; a \neq b \to \neg(a \leq b) \wedge \neg(b \leq a)$.
\end{remark*}

\begin{definitionbox}{Definition (Initial Segment)}
\begin{definition}[Initial Segment]\label{def:initial-segment}
A subset $I \subseteq A$ is an \emph{initial segment} of $(A,\leq)$ if
\[
a \in I \text{ and } b \leq a \;\Rightarrow\; b \in I.
\]
\end{definition}
\end{definitionbox}

\NoLocalDependencies

\begin{remark*}[Standard quantified statement]
$I$ is an initial segment $\;\leftrightarrow\; \forall a \in I,\; \forall b \in A,\; b \leq a \to b \in I$.
\end{remark*}

\begin{remark*}[Transition]
Ordered sets provide the abstract framework for order structures on the real
numbers, function spaces, and metric spaces. In later sections, order
interacts with topology and analysis through intervals, monotone functions,
and the completeness axiom for $\mathbb{R}$.
\end{remark*}

% ---------------------------------------------------------
% Directed Sets and Nets
% ---------------------------------------------------------

\begin{definitionbox}{Definition (Directed Set)}
\begin{definition}[Directed Set]\label{def:directed-set}
A preorder $(I, \leq)$ is \emph{directed} (or a \emph{directed set}) if every
finite subset of $I$ has an upper bound in $I$: for all $i, j \in I$ there
exists $k \in I$ with $i \leq k$ and $j \leq k$.
\end{definition}
\end{definitionbox}

\begin{remark*}[What directed means]
In a directed set, any two elements can always be ``dominated'' by a third.
The condition does not require a maximum element to exist, only that no two
elements are ever incomparable in a final way: there is always a common
upper bound available. Every totally ordered set is directed (take $k =
\max(i,j)$), but directed sets can be far more general.
\end{remark*}

\begin{example*}[Canonical directed sets]
\begin{enumerate}[label=(\roman*)]
\item $(\mathbb{N}, \leq)$ is directed: given $m, n \in \mathbb{N}$, take
$k = \max(m,n)$.
\item The collection of all finite subsets of a set $X$, ordered by inclusion,
is directed: given finite sets $F_1, F_2 \subseteq X$, take $k = F_1 \cup
F_2$.
\item The set of all open neighborhoods of a point $x$ in a topological space,
ordered by reverse inclusion ($U \leq V$ iff $U \supseteq V$), is directed:
given neighborhoods $U$ and $V$ of $x$, take $k = U \cap V$, which is also
a neighborhood of $x$.
\item A partially ordered set is directed if and only if every finite subset
has an upper bound. Partially ordered sets that are \emph{not} directed have
``incomparable maximal elements'' — elements with no common upper bound.
\end{enumerate}
\end{example*}

\begin{definitionbox}{Definition (Net)}
\begin{definition}[Net]\label{def:net}
Let $(I, \leq)$ be a directed set and $X$ a set. A \emph{net} in $X$ is a
function $x : I \to X$. The value at $i \in I$ is written $x_i$, and the
net is written $(x_i)_{i \in I}$.
\end{definition}
\end{definitionbox}

\begin{remark*}[Nets generalize sequences]
A sequence is a net with index set $(\mathbb{N}, \leq)$. Nets are the correct
generalization of sequences to general topological spaces: in a metric space,
sequences suffice to detect all topological properties (closure, continuity,
compactness). In a general topological space, sequences can fail to detect
these properties — one needs nets. Specifically, a subset $A$ of a
topological space is closed if and only if every convergent net in $A$ has
its limit in $A$; and a function is continuous if and only if it preserves
limits of nets.

The definition of net is placed here because it is purely order-theoretic:
a directed index set and a function. The analytic content — what it means
for a net to converge — belongs to topology and functional analysis. By
establishing the set-theoretic definition now, those chapters can immediately
state convergence conditions without set-theoretic digression.
\end{remark*}
